\section{Chapter 3 State of the Art and Existing Cloud Security Solutions}

\subsection{3.1 Introduction}

The rapid adoption of cloud computing has fundamentally transformed organizational information technology infrastructures. As enterprises increasingly migrate critical workloads to public and hybrid cloud environments, ensuring the security of cloud resources has become one of the most significant challenges in modern cybersecurity. Cloud infrastructures consist of numerous interconnected services, identities, storage systems, virtual networks, databases, application workloads, and serverless resources, each containing hundreds of configurable security parameters. To address these challenges, cloud providers and cybersecurity vendors have developed various Cloud Security Posture Management (CSPM) and Cloud-Native Application Protection Platform (CNAPP) solutions capable of continuously monitoring cloud environments for security misconfigurations and compliance violations. These platforms assist organizations by automatically discovering cloud assets, evaluating security configurations, enforcing compliance standards, and generating remediation recommendations. Although existing solutions significantly improve cloud visibility, most continue to rely primarily on rule-based security assessment methodologies. Consequently, they often generate extensive lists of isolated findings without adequately considering relationships among cloud resources or evaluating how multiple configuration weaknesses collectively contribute to exploitable attack paths. This chapter presents a detailed analysis of widely adopted cloud security solutions, examining their architecture, capabilities, strengths, and limitations. Understanding the current state of cloud security technology provides the foundation for identifying the research gap addressed by the proposed CloudSentinel AI framework.

\subsection{3.2 Amazon Web Services Security Hub}

AWS Security Hub is Amazon's centralized cloud security management service designed to provide a unified view of security findings across AWS accounts and regions. It aggregates security findings generated by various AWS security services, including Amazon GuardDuty, Amazon Inspector, AWS IAM Access Analyzer, AWS Firewall Manager, and AWS Macie. Security Hub continuously evaluates cloud resources against established security frameworks such as the AWS Foundational Security Best Practices, CIS AWS Foundations Benchmark, PCI DSS, and NIST standards. Identified security findings are prioritized according to predefined severity levels and presented through a centralized dashboard. The primary objective of Security Hub is to simplify cloud security monitoring by consolidating findings from multiple AWS security services into a single management interface.\\
\textbf{Key Features}

\begin{itemize}
\tightlist
\item
  Centralized security dashboard\\
\item
  Continuous compliance monitoring\\
\item
  Multi-account security aggregation\\
\item
  Integration with AWS security services\\
\item
  Security standards assessment\\
\item
  Automated security findings
\end{itemize}

\textbf{Advantages}

\begin{itemize}
\tightlist
\item
  Native integration with AWS services.\\
\item
  Easy deployment within AWS environments.\\
\item
  Comprehensive compliance reporting.\\
\item
  Continuous monitoring of AWS resources.\\
\item
  Automated collection of security findings.
\end{itemize}

\textbf{Limitations}

\begin{itemize}
\tightlist
\item
  Primarily focused on AWS environments.\\
\item
  Relies heavily on predefined security rules.\\
\item
  Generates large numbers of independent alerts.\\
\item
  Limited contextual understanding of cloud resources.\\
\item
  Does not construct complete attack paths across cloud assets.\\
\item
  Limited AI-assisted explanation capabilities.
\end{itemize}

\subsection{3.3 AWS Config}

AWS Config is a configuration management service that continuously records the configuration state of AWS resources and monitors configuration changes over time. It enables organizations to evaluate cloud resources against predefined configuration rules and maintain historical records for auditing and compliance purposes. AWS Config provides continuous configuration recording, change tracking, and compliance evaluation. Security administrators can define custom rules to detect insecure configurations, monitor resource drift, and enforce organizational security policies.\\
\textbf{Key Features}

\begin{itemize}
\tightlist
\item
  Configuration history\\
\item
  Resource inventory\\
\item
  Configuration drift detection\\
\item
  Compliance evaluation\\
\item
  Custom rule creation\\
\item
  Integration with AWS Lambda
\end{itemize}

\textbf{Advantages}

\begin{itemize}
\tightlist
\item
  Excellent configuration auditing.\\
\item
  Historical tracking of configuration changes.\\
\item
  Supports automated compliance monitoring.\\
\item
  Enables custom policy development.\\
\item
  Useful for regulatory compliance.
\end{itemize}

\textbf{Limitations}

\begin{itemize}
\tightlist
\item
  Focuses primarily on configuration state.\\
\item
  Limited contextual risk analysis.\\
\item
  Does not correlate multiple security findings.\\
\item
  No attack graph generation.\\
\item
  Minimal prioritization beyond rule evaluation.
\end{itemize}

\subsection{3.4 Prowler}

Prowler is an open-source cloud security assessment tool widely used for evaluating AWS security configurations against industry-recognized security benchmarks. It performs automated security assessments based on CIS Benchmarks, AWS Well-Architected Framework recommendations, PCI DSS, ISO 27001, HIPAA, and numerous additional compliance standards. Prowler executes hundreds of individual security checks covering identity management, networking, encryption, storage, monitoring, logging, and cloud governance.\\
\textbf{Key Features}

\begin{itemize}
\tightlist
\item
  Open-source\\
\item
  CIS Benchmark support\\
\item
  Multi-framework compliance\\
\item
  Command-line execution\\
\item
  HTML and CSV reporting\\
\item
  Automated security scanning
\end{itemize}

\textbf{Advantages}

\begin{itemize}
\tightlist
\item
  Free and open source.\\
\item
  Large collection of security checks.\\
\item
  Supports automation.\\
\item
  Easy integration into CI/CD pipelines.\\
\item
  Community-driven development.
\end{itemize}

\textbf{Limitations}

\begin{itemize}
\tightlist
\item
  Rule-based detection only.\\
\item
  Produces static reports.\\
\item
  Limited visualization.\\
\item
  No contextual attack analysis.\\
\item
  No AI-assisted prioritization.
\end{itemize}

\subsection{3.5 ScoutSuite}

ScoutSuite is an open-source multi-cloud security auditing platform that provides read-only security assessments across Amazon Web Services (AWS), Microsoft Azure, Google Cloud Platform (GCP), Oracle Cloud Infrastructure (OCI), and Kubernetes environments. Unlike provider-specific tools, ScoutSuite offers a unified security assessment interface across multiple cloud providers, making it particularly useful for organizations operating hybrid or multi-cloud infrastructures. ScoutSuite generates an interactive HTML report that enables administrators to explore cloud resources and identify security weaknesses.\\
\textbf{Key Features}

\begin{itemize}
\tightlist
\item
  Multi-cloud support\\
\item
  Interactive reports\\
\item
  Read-only security assessment\\
\item
  IAM analysis\\
\item
  Network analysis\\
\item
  Storage assessment
\end{itemize}

\textbf{Advantages}

\begin{itemize}
\tightlist
\item
  Supports multiple cloud providers.\\
\item
  Interactive visualization.\\
\item
  Easy deployment.\\
\item
  Comprehensive inventory collection.\\
\item
  Useful for security audits.
\end{itemize}

\textbf{Limitations}

\begin{itemize}
\tightlist
\item
  Static reporting.\\
\item
  Limited contextual reasoning.\\
\item
  No attack graph generation.\\
\item
  Limited prioritization.\\
\item
  Does not model cloud relationships.
\end{itemize}

\subsection{3.6 Wiz}

Wiz is a modern Cloud-Native Application Protection Platform (CNAPP) that provides comprehensive cloud security visibility through agentless scanning, graph-based asset relationships, vulnerability management, identity analysis, and attack path visualization. Unlike traditional CSPM platforms, Wiz correlates cloud assets, vulnerabilities, permissions, and network exposure to identify toxic combinations that significantly increase organizational risk. Its Security Graph architecture enables visualization of relationships among cloud resources, providing security analysts with a broader understanding of cloud attack surfaces.\\
\textbf{Key Features}

\begin{itemize}
\tightlist
\item
  Agentless architecture\\
\item
  Security Graph\\
\item
  Attack path visualization\\
\item
  Vulnerability management\\
\item
  Identity analysis\\
\item
  Multi-cloud support
\end{itemize}

\textbf{Advantages}

\begin{itemize}
\tightlist
\item
  Strong visualization.\\
\item
  Graph-based architecture.\\
\item
  Excellent cloud inventory.\\
\item
  Modern CNAPP capabilities.\\
\item
  Good prioritization.
\end{itemize}

\textbf{Limitations}

\begin{itemize}
\tightlist
\item
  Commercial platform with licensing costs.\\
\item
  Limited transparency regarding proprietary risk-scoring algorithms.\\
\item
  AI explanation capabilities remain relatively limited.\\
\item
  Graph analysis is primarily optimized for platform-specific workflows.\\
\item
  Custom research extensions are difficult.
\end{itemize}

\subsection{3.7 Prisma Cloud}

Prisma Cloud, developed by Palo Alto Networks, is one of the most comprehensive enterprise CNAPP platforms available today. It integrates CSPM, Cloud Workload Protection (CWPP), Infrastructure-as-Code scanning, container security, Kubernetes security, API security, and identity management into a unified cloud security platform. Prisma Cloud provides extensive compliance monitoring, vulnerability assessment, runtime protection, and cloud governance capabilities.\\
\textbf{Advantages}

\begin{itemize}
\tightlist
\item
  Comprehensive security coverage.\\
\item
  Enterprise-grade scalability.\\
\item
  Strong compliance support.\\
\item
  Excellent workload protection.\\
\item
  Multi-cloud architecture.
\end{itemize}

\textbf{Limitations}

\begin{itemize}
\tightlist
\item
  High deployment complexity.\\
\item
  Commercial licensing costs.\\
\item
  Large volume of security findings.\\
\item
  Limited explainability.\\
\item
  Context remains constrained by predefined analytical models.
\end{itemize}

\subsection{3.8 Microsoft Defender for Cloud}

Microsoft Defender for Cloud provides integrated cloud security management for Microsoft Azure while also supporting AWS and Google Cloud environments. The platform combines CSPM, workload protection, vulnerability assessment, and compliance management. It continuously evaluates cloud resources using Microsoft Secure Score and provides remediation recommendations based on organizational security posture.\\
\textbf{Advantages}

\begin{itemize}
\tightlist
\item
  Native Azure integration.\\
\item
  Secure Score.\\
\item
  Regulatory compliance support.\\
\item
  Threat protection.\\
\item
  Hybrid cloud support.
\end{itemize}

\textbf{Limitations}

\begin{itemize}
\tightlist
\item
  Optimized primarily for Azure.\\
\item
  Secure Score remains largely rule-based.\\
\item
  Limited attack path reasoning.\\
\item
  AI explanation capabilities are basic.
\end{itemize}

\subsection{3.9 Comparative Analysis of Existing Solutions}

The analysis of current cloud security platforms demonstrates that significant progress has been made in automating cloud security assessments. Most solutions effectively identify cloud misconfigurations, enforce compliance standards, and improve visibility across cloud infrastructures. However, they share several common limitations that motivate further research. First, the majority of platforms evaluate cloud resources independently rather than reasoning about the relationships among assets. Consequently, security findings are often presented as isolated alerts, making it difficult for analysts to understand how multiple weaknesses interact. Second, most platforms prioritize findings using predefined severity rules instead of contextual risk analysis. As a result, organizations may spend valuable time remediating low-impact findings while overlooking attack chains that pose substantially greater risk. Third, although graph-based approaches have been introduced in some commercial platforms, their methodologies are generally proprietary and provide limited transparency regarding risk calculation and decision-making. Finally, current platforms offer only limited explainability. Security analysts frequently receive alerts describing configuration violations without sufficient context regarding the attack path, business impact, or rationale behind the assigned severity. These limitations indicate that there remains an opportunity to develop more intelligent cloud security frameworks capable of integrating contextual reasoning, graph-based analysis, dynamic risk assessment, and AI-assisted explanation.

\subsection{3.10 Chapter Summary}

This chapter examined the current state of cloud security technologies by analyzing leading commercial and open-source cloud security platforms. The evaluation demonstrates that while existing CSPM and CNAPP solutions have significantly improved cloud security monitoring, they continue to exhibit limitations in contextual reasoning, attack path generation, dynamic risk prioritization, and explainable security analysis. These observations provide the foundation for the next chapter, which critically reviews existing academic research and identifies the specific research gap that the proposed CloudSentinel AI framework seeks to address.
